\documentclass[bachelor,paper]{pmk-graduate}


\usepackage{amsthm}
\usepackage{amsmath} 
\usepackage[shortcuts]{extdash}
\usepackage{graphicx}
\usepackage{float}
\usepackage{mathtext}
\usepackage{amssymb}
\usepackage{algorithm,algorithmicx,algpseudocode}

\addbibresource{milov.bib}

\newtheorem{thm}{Теорема}
\theoremstyle{definition}
\newtheorem{dfn}{Определение}

%\counterwithout{figure}{section}
%\counterwithout{figure}{chapter}

% \toggletrue{bbx:gostbibliography}

\begin{document}
% Учебный год
\year{2025-2026}

% Приказ: дата, номер
%\order{02.12.2025}{222-С}

% Дисциплина для курсовой работы
%\subject{Дискретная математика}

% Студент
\student{Милов Данила Константинович
\brief Д.\,К.\,Милов
\group 47
\english Milov Danila Konstantinovich
}

% Научный руководитель
\superviser{Дудаков Сергей Михайлович
\brief С.\,М.\,Дудаков
\position декан факультета ПМиК
\degree д.\,ф.-м.\,н. % Учёная степень. При отсутствии убрать строку
\rank профессор % Учёное звание. При отсутствии убрать строку
}

% Образовательная программа. См. перечень ниже
\program{ics}
% Бакалавриат:
% ics - Информатика и компьютерные науки
% swi - Инженерия программного обеспечения
% swiai - Программная инженерия в искусственном интеллекте
% aida - Искусственный интеллект и анализ данных
% mm - Математическое моделирование
% sa - Системный анализ
% icpim - Инженерное компьютерное проектирование и индустриальная математика
% acse - Прикладная информатика в экономике
% acsm - Прикладная информатика в мехатронике
% icmrs - Интеллектуальное управление в мехатронных и робототехнических системах
% Магистратура:
% macsae - Прикладная информатика в аналитической экономике
% mista - Интеллектуальные системы. Теория и приложения
% msp - Системное программирование
% msa - Системный анализ
% mitcd - Информационные технологии в управлении и принятии решений


% Заголовок ВКР
\title{Изучение методов фрактального сжатия для различных типов информации
    \english Fractal compression of different types of information}

% Цель ВКР
\goal{Изучить и реализовать алгоритмы фрактального сжатия для различных типов информации.}

% Задачи для разработки в ВКР
\begin{tasks}
  \task{Изучение алгоритмов фрактального сжатия для изображений.
  \result
  Изучены алгоритмы сжатия изображений. Были рассмотрены классические алгоритмы сжатия, такие как PNG или JPEG, а также предмет изучения данной работы — алгоритм фрактального сжатия при помощи PIFS, и его модификации.
  }
  \task{Изучение алгоритмов фрактального сжатия звука.
  \result
  Изучены алгоритмы сжатия аудиоданных. Кратко рассмотрен принцип работы сжатия MP3, FLAC, кроме того была исследована тема применимости фрактального сжатия к аудиоданным и особенности этого алгоритма.
  }
  \task{Реализация алгоритмов фрактального сжатия.
  \result
  На основе изученных алгоритмов фрактального сжатия реализованы программы для компрессии этих данных. Программа fcomp-audio для сжатия звука принимает на вход файлы в формате WAV, а fcomp-img работает с JPEG и PNG форматами. Приложения реализуют алгоритм с разделением на блоки фиксированного размера — один из простейших вариантов разбиения в PIFS, который, тем не менее, даёт вполне приемлемые результаты.
  }
  \task{Сравнение качества сжимающих алгоритмов.
  \result
  Было проведено сравнение реализованных алгоритмов, а также проведены измерения степени сжатия и количества потерь при работе программы.
  }
  \task{Адаптация алгоритмов фрактального сжатия для других типов информации.
  \result
  Исследован вопрос применимости фрактального кодирования к другим предметным областям. Рассмотрена абстрактная версия алгоритма фрактального сжатия, описаны вопросы, на которые необходимо ответить при разработке версии алгоритма для какой-то конкретной предметной области, а также описаны некоторые исследования применимости указанного алгоритма к другим предметным областям.
  }
\end{tasks}

% Перечень иллюстративного материала
%\begin{images}
%\image{1}    
%\image{2}    
%\end{images}

% Консультанты
%\begin{advisers}
%\adviser{Сидоров Сидор Сидорович, проектирование схемы БД}    
%\end{advisers}

% Команда для генерации задания на ВКР
%\makeassignment
%\end{document}

\begin{englishabstract}
This work is focused on finding information about different algorithms for compression of different types of data, such as images and audiofiles. Mainly, we talk about fractal compression algorithm. In this paper we look at classic fractal method of compressing information and some of it's modifications that can improve compression ratio, file sizes or decrease the loss of quality of the resulting image. We provide a reader with mathematical bckground, needed for understanding the algorithm, including some crucial theorems and definitions from compressive mappings and fractal theory.
  Besides, we make our own implementation of PIFS-based fractal encoding method for both images and audio files, that is quite primitive and is more of a demonstration of the idea, but powerfull nonetheless. Other than that, we also create our own binary file format for storing fractal codes in a more compact manner than it's plain text form.
  In this paper we thoroughly describe each step of the algorithm, discuss it both without and with relation to images and audio, trying to show that it's idea can be used in various domains and areas. We also look at the modern state of using fractal data compression algorithm in other fields than audio and images.
\end{englishabstract}

\renewcommand{\thealgorithm}{\arabic{algorithm}.}

\maketitlepage
\introduction
\actuality
Проблема хранения и передачи информации является одной из самых важных с самого зарождения информатики. Поскольку файлы в исходном виде занимают очень много места, критически необходимо применять методы, способные сократить объём таких файлов. Такие методы называются сжатием --- процессом уменьшения занимаемого данными места путём уменьшения избыточности информации. А поскольку практически все типы информации обладают избыточностью в том или ином виде: пиксели совпадающих цветов на изображении, повторяющиеся буквы и слова в тексте, совпадающие фрагменты аудиоданных, сжатие является крайне эффективным способом экономии места на диске или пропускной способности соединения.

Существует огромное количество алгоритмов сжатия информации, напри-
мер, PNG и JPEG для фото, или MP3 и FLAC для аудио. Все они являются широко используемыми и хорошо зарекомендовавшими себя алгоритмами. Помимо них существует также гораздо менее известный алгоритм, который не получил большой популярности в годы открытия по причине необходимости больших вычислительных мощностей, недоступных в то время. Однако сейчас, с ростом производительности ЭВМ и исследованиями, направленными на улучшение алгоритма, он является актуальным и перспективным способом сжатия информации. Речь идёт о фрактальном сжатии информации, идеи которого впервые были предложены в работах Майкла Барнсли [1, 2].

Этот вид компрессии обладает рядом интересных свойств и потенциально может позволить достичь высоких степеней сжатия при сохранении приемлемого качества восстановленных данных. Среди достоинств можно выделить следующие.
\begin{itemize}
  \item Возможность манипулировать разрешением разархивированного изображения/аудиофайла(получать изображение/звук более или менее детализированный, чем оригинал).
  \item Быстрый процесс декодирования.
  \item Потенциально большие степени сжатия с малыми потерями качества.
\end{itemize}


\review
В ходе исследования теории, необходимой для выполнения данной работы, были изучены работы Майкла Барнсли, Джона Хатчинсона и Арно Жакена, стоявших у истоков развития фрактального сжатия информации [1, 2, 7]. Поскольку наиболее широкое распространение получил именно предложенный Арно Жакеном способ нахождения сжимающего отображения, кодирующий изображение, было решено имплементировать именно этот вариант алгоритма. Для лучшего понимания математики, стоящей за процессами сжатия и декодирования, было принято решение сначала изучить теорию сжимающих отображений [3], а также теорию фракталов [4, 5]. После более глубокого изучения теории, был произведён поиск материалов, связанных с оптимизацией и модификацией алгоритма [8, 9]. Для программной реализации был выбран язык программирования Go, поэтому очень важными ресурсами были официальная документация языка [13], интерактивный курс по возможностям, особенности и структуре языка [14], а также сайт с информацией и документацией по всем доступным пакетам языка Go [10]. Также в рамках исследования применимости алгоритма к предметным областям, алгоритмы для которых не были представлены в тексте настоящей работы, были просмотрены статьи [11, 12], показывающие возможность применимости фрактального сжатия не только к изображениям и звуку.

\goalandtasks

\tools
В ходе выполнения работы были изучены различные вариации алгоритма фрактального сжатия для изображений и звука, а также их модификации. В качестве средств создания программной реализации был выбран язык программирования Go (также известный, как Golang), его богатая стандартная библиотека (в особенности пакет image), а также некоторые сторонние пакеты, такие как go-flag для создания cli-интерфейса, go-wav для работы с форматом аудиоданных WAV. Помимо этого для конвертации из MP3 и других форматов в WAV использовалась программа FFmpeg, свободная, открытая и доступная на всех операционных системах.

\results

Кроме того, в ходе работы для программ написан удобный cli-интерфейс, позволяющий управлять параметрами алгоритма. Он позволяет настроить размеры блоков, указать входные и выходные файлы, а также установить пороговое значение для перебора и выделить количество потоков, используемых программами. Также был разработан собственный формат файлов, позволяющий компактно записывать сжатые данные. Информация в нём хранится в бинарном представлении, что позволяет сильно сэкономить в объёме файла, если сравнивать с текстовым представлением. Помимо этого, разработанный формат пользуется особенностями алгоритма, чтобы не записывать явно информацию, которую можно после восстановить аналитически.

\structure
Настоящая работа состоит из введения, двух крупных глав, посвящённых: теории и основным понятиям, а также описанию программного обеспечения. Введение кратко описывает актуальность исследований в рассматриваемой сфере, комментирует то, как были использованы различные литературные источники, перечисляет цели и задачи исследования, кратко комментирует методы и средства, которыми мы воспользовались, а также освещает основные результаты, достигнутые в ходе работы. Глава «Теория и основные определения» посвящена рассмотрению определений и теорем, критически важных для понимания изложенного далее материала, а также рассказывает общую суть алгоритма фрактального сжатия. Вторая глава, в свою очередь, посвящена результатам работы в рамках данного исследования. Сначала в ней рассказывается о том, каким планировалось увидеть итоговый вариант программ сжатия, затем речь идёт о выборе средств разработки и причинах, по которым было решено остановиться именно на Go. Следующим блошим разделом второй главы является описание алгоритма фрактального сжатия уже в контексте текущей работы, в этом разделе сначала даётся общая структура алгоритма, а затем более подробно рассматриваетя каждый этап и описываются решения, на которых было решено остановиться при создании программ для компрессии. Далее следует часть, посвящённая внутренней структуре программ на Go и описанию того, какие части алгоритма воплощают те или иные функции в коде. После чего кратко даётся комментарий о применимости фрактального сжатия к другим предметным областям и текущем состоянии исследований по этой теме. Наконец, завершает вторую главу демонстрация результатов работы и тесты потерь качества, размеров файлов и времени работы.

Также в работу включено два приложения. В приложении А можно ознакомиться с кодом программ, реализующих описанные далее алгоритмы, а приложение В демонстрирует интерфейс этих программ, позволяющий удобно управлять ходом работы программ.

\chapter{Теория и основные определения}
В этой главе представлена теоритическая основа алгоритмов фрактального сжатия в отрыве от традиционных способов компрессии. Здесь мы опишем основную идею, математические конструкции и теоремы, связанные с тематикой работы.

\section{Сжимающие отображения}
\begin{dfn}[Сжимающее отображение \cite{funcAn}]
  Пусть $(E,d)$ --- полное метрическое пространство, а $f:E\rightarrow E$ --- отображение $E$ на $E$. Мы говорим, что $f$ является \emph{сжимающим отображением}, если существует $0 < \alpha < 1$, такое, что:
  $$\forall x,y \in E, \rho(f(x), f(y)) \leqslant \alpha \rho(x,y).$$
\end{dfn}

О сжимающих отображениях существует две важные в рамках данной работы теоремы: теорема Банаха о неподвижной точке и теорема коллажа.
\begin{thm}[Теорема Банаха \cite{funcAn}]
  Пусть $(E,d)$ --- полное метрическое пространство, а $f:E\rightarrow E$ --- сжимающее отображение над $E$. Тогда существует единственная неподвижная точка отображения $f$.
\end{thm}
\begin{thm}[Теорема коллажа]
  Пусть $(E,d)$ --- полное метрическое пространство. $f$ --- сжимающее отображение над $E$. Тогда: $$\text{если } \rho(x,f(x))<\varepsilon, \text{ то } \rho(x,x_0) < \dfrac{\varepsilon}{1-s}.$$
\end{thm}

Теорема коллажа говорит нам, что если мы найдём такое сжимающее отображение $f$, что $f(x)$ достаточно близко к $x$, то мы можем быть уверены, что неподвижная точка $f$ тоже находится близко к $x$.

Этот результат и будет основой для всей дальнейшей работы, поскольку предоставляет способ нахождения отображения, чья неподвижная точка близка к сжимаемой информации. Суть кодирования в алгоритме в следующем: вместо того, чтобы сохранять сами данные, мы можем сохранить только такое сжимающее отображение, по которому можно восстановить то, что было сжато.

\section{Фракталы}
Прежде чем переходить к реализации, необходимо дать несколько определений и некоторую историческую информацию, чтобы понять, почему алгоритм получил такое название, при чём здесь фракталы и как он был открыт. Начнём с определения фрактала.

\begin{dfn}[Фрактал \cite{fractalDef}]
  Множество, обладающее свойством самоподобия или масштабной инвариантности (объект в точности или приближённо совпадающий с частью себя самого). Масштабно-инвариантные системы обычно характеризуются нецелой --- {\em фрактальной размерностью}.
\end{dfn}

Фракталы изначально не были придуманы с целью сжимать данные. Понятие нецелой размерности, а также некоторые свойства фрактальных объектов изучались такими выдающимися учёными, как Георг Кантор, Джузеппе Пеано, Давид Гильберт, Феликс Хаусдорф и многими другими. Сам же термин ввёл Бенуа Мандельброт в семидесятых годах прошлого века \cite{fractalDef, mandelbrot}.

Далее определим важное для понимания алгоритма понятие системы итерируемых функций (СИФ, англ. Iterated Function Systems или IFS).
\begin{dfn}[Система итерируемых функций \cite{fractalDef}]
  Это система функций из некоторого фиксированного класса функций, отображающих одно многомерное множество на другое.
\end{dfn}

Идея использования СИФ для представления фракталов впервые была представлена в статье Джона Хатчинсона <<Фракталы и самоподобие>> \cite{hutchinson}.%https://maths-people.anu.edu.au/~john/Assets/Research%20Papers/fractals_self-similarity.pdf


Большую роль в исследовании СИФ и их применении для сжатия информации сыграли исследования Майкла Барнсли, который в своей статье в журнале Byte за январь 1988 г. \cite{barnsleyByte} предложил метод для сжатия изображений путём представления их в виде некоторой системы итерируемых функций. Он доказал, что СИФ может быть использована для кодирования изображения, которое после можно восстановить благодаря теореме Банаха о неподвижной точке. То есть он показал, что существуют такие системы функций, для которых некоторое изображение является неподвижной точкой, а значит при любом начальном значении в результате многократного применения этих функций мы получим закодированное изображение. Однако в этой статье не было описано метода построения этой самой СИФ для произвольных изображений.

Способ решения данной проблемы был предложен учеником Майкла Барнсли, Арно Жакеном \cite{jacquin}. Он описал метод, по которому можно использовать не преобразования всего изображения, а сопоставлять отдельные его части. Предложенный метод получил название <<Сис\-те\-ма итерируемых ку\-соч\-но-оп\-ре\-де\-лён\-ных функ\-ций>> (Partitioned Iterated Function System --- PIFS).

\section{PIFS и алгоритм фрактального сжатия}
Использование PIFS позволяет получать искомые сжимающие отображения для произвольных данных. Далее будет описан общий алгоритм действий.
\begin{enumerate}
  \item Данные разбиваются на непересекающиеся полностью покрывающие их блоки $R_1,\dots,R_n$. Эти блоки называются ранговыми (от англ. range --- область значений функции) или конечными, интервальными.
  \item Далее те же данные разбиваются на блоки источников или доменные блоки $D_1,\dots, D_m$. Они необязательно разделены, необязательно покрывают все данные.
  \item После этого необходимо для каждого $R_i$ выбрать $D_{j_i}$ и отображение $f_i: [0,1]^{D_{j_i}}\rightarrow[0,1]^{R_i}$.
  \item Наконец, мы можем определить нашу функцию $f$ как $f(x)_k = f_i(x_{D_{j_i}})_k$ для $k\in R_i$.
\end{enumerate}

Выдвигается следующее утверждение: если все $f_i$ являются сжимающими отображениями, то и $f$ тоже сжимающее отображение. Это следует из того, что $R_i$ --- попарно непересекающиеся блоки. Полное доказательство опускается.

Перед нами остаётся один вопрос, который нужно рассмотреть: как выбрать $D_i$ и $f_i$? Для этого следует обратиться к ранее упомянутой теореме коллажа, предоставляющей способ их выбора: если $x_{R_i}$ достаточно близко к $f_i(x_{D_i})$ для всех $i$, то $x$ находится близко к $f(x)$, а значит по теореме коллажа $x$ близок к $x_0$.

Получается, что можно независимо для каждого $R_i$ построить множество сжимающих отображений $D_i$ на него и выбрать наилучшее. В этом и состоит идея алгоритма. Далее будут более подробно описаны шаги алгоритма сжатия и декодирования в контексте тех типов информации, для которых они разработаны.

\chapter{Описание программного обеспечения}

В ходе работы были написаны программы для сжатия изображений и звука, реализующие алгоритм фрактальной компрессии при помощи PIFS, изложенный выше.

\section{Требования, предъявляемые к программному обеспечению}

В ходе планирования и изучения материалов были сформулированы следующие требования, в соответствии с которыми велась разработка.
\begin{itemize}
  \item Программы должны реализовывать алгоритм фрактального сжатия.
  \item Размер сжатых данных должен быть сравним хотя бы с алгоритмами сжатия без потерь (быть не сильно хуже или даже лучше).
  \item Приемлемое качество восстановленной информации.
  \item Возможность отображения подсчитанных коэфициентов преобразования в виде JSON-файла для отладки и демонстрации работы алгоритма.
  \item Необходим удобный консольный интерфейс с возможностью настройки параметров сжатия.
\end{itemize}

\section{Выбор средств разработки}
Алгоритм фрактального сжатия подразумевает огромное количество перебора, но при этом отлично распараллеливается. Поскольку задача компрессии крайне ресурсозатратная, вопроса об использовании каких-либо языков, кроме компилируемых, не стояло. Но среди них уже был большой выбор: есть хорошо зарекомендовавшие себя в ресурсоёмких задачах C и C++, есть другие перспективные языки, такие как Rust и Go. В итоге выбор пал на Go \cite{godoc}. Среди его преимуществ можно выделить нижеследующие.
\begin{itemize}
  \item Богатая стандартная библиотека. Например, пакет для работы с изображениями, лежащий в основе одной из частей этого проекта, написан при помощи этого пакета.
  \item Удобная работа с многопоточностью благодаря высоко- и низкоуровневым инструментам параллелизации.
  \item Наличие сборщика мусора (garbage collector). Это не было основной причиной, однако это добавляет удобства в некоторых ситуациях. На производительности сборщик мусора оставляет практически незаметный эффект.
  \item Наличие установщика пакетов и широкой библиотеки модулей для всевозможных задач. Это было одной из главных причин, поскольку возможность скачать любой модуль, просто указав ссылку на github в своём файле --- это очень удобно.
  \item Простой, минималистичный и понятный синтаксис, а также низкий порог входа и быстрое обучение.
\end{itemize}

Помимо этого, при работе с аудиоданными использовалась программа FFmpeg для конвертации из различных форматов звука в WAV, с которым может работать приложение.

\section{Вариант алгоритма, реализованный в работе}

Для начала стоит определить составные части алгоритма фрактального сжатия. Все они допускают различные реализации и возможность модификации. Ниже представлен краткий список основных элементов, а далее мы подробно остановимся на каждом из них в контексте той или иной предметной области.
\begin{itemize}
  \item Метод разбиения. Диктует то, как информация разбивается на ранговые и доменные блоки.
  \item Метрика, используемая для проверки соответствия рангового блока и доменного, к которому применили преобразование.
  \item Преобразования. Описывают различные изменения фрагментов, позволяющие добиться более близкого отображения.
  \item Метод перебора. Поскольку перебор блоков --- наиболее времязатратная часть процесса сжатия, ведутся исследования по его оптимизации.
  \item Декодирование. Восстановление исходного изображения из фрактальных кодов.
  \item Хранение коэфициентов преобразований в файле. Помимо нахождения фрактальных кодов, нужно также компактно записывать их и считывать.
\end{itemize}


%Написанная программа работает с изображениями формата PNG и JPEG, сжимает их при помощи предложенного выше алгоритма и выдаёт файл, содержащий информацию о преобразовании либо в виде JSON-документа, либо в придуманном в ходе разработки бинарном формате. Поддерживаются различные настройки параметров алгоритма, такие как: размеры блоков, названия входных и выходных файлов, количество потоков, используемых для перебора, а также возможность указать пороговое значение различий между ранговым и отображением доменного, при котором останавливается дальнейший перебор.

\subsection{Метод разбиения}

Начнём с решения вопоса о методе разбиения. От выбора того, как выделять доменные и ранговые блоки, зависит качество сжатия, время работы, объёмы перебора и количество соотношений в результате.

Сущетвует несколько стратегий, применимых в сжатии изображений и звука при помощи PIFS. Самым простым является разбиение на блоки фиксированного размера. Это и есть способ, на котором я решил остановиться. Он позволяет добиться приемлемых результатов, при этом очень прост в программировании, а также оказался хорошим выбором для создания компактного бинарного файла, потому что размеры блоков можно записать только один раз без необходимости уточнять размеры блоков в каждом конкретном преобразовании. В изображении выделяются одинаковые квадратные сегменты, а в случае аудиофайлов деление идёт на отрезки равной длины.

\begin{figure}[H]
  \begin{center}
    \includegraphics[width=0.65\textwidth]{./images/PIFS.png}
  \end{center}
  \caption{Разбиение изображения на блоки фиксированного размера}
\end{figure}

Помимо этого популярным решением для изображений является использование подхода с разбиением при помощи квадродерева. При таком подходе изображение может быть разбито на клетки разного размера. Для реализации такого метода необходимо задать параметр точности совпадения $\varepsilon$, с которым будет сравниваться метрика на итерациях. В таком случае алгоритм разбиения следующий.
\begin{enumerate}
  \item Разбить изображение на доменные блоки любым способом.
  \item Задать начальное разбиение (например, ранговые блоки 64$\times$64).
  \item Если после перебора не нашлось такого преобразования доменного блока в ранговый блок $R_i$, что значение метрики меньше, чем $\varepsilon$, то $R_i$ разбивается на 4 части: $R_{NW}, R_{NE}, R_{SE}, R_{WE}$, для каждой из которых производится перебор доменных блоков и преобразований.
\end{enumerate}

Алгоритм с использванием квадродерева может иметь более длительное время работы в худшем случае или более быстрое в лучшем, однако его главным преимуществом является то, что фрагменты изображения, обладающие малым числом деталей (например, однотонный фон), кодируются экономно, а фрагменты, содержащие большое количество визуальной информации, получают наибольшее внимание, позволяя в итоге получить компромисс между качеством изображения и объёмом сжатого файла. Использование квадродерева и некоторых оптимизаций могут привести к заметной экономии места \cite{quadTreeOptimiz}. А поскольку реализация, изложенная в настоящей работе, уже имеет возможность ввода порогового значения для оптимизации перебора, модификация разбиения с использованием квадродерева --- логичное развитие и хорошая идея для будущих исследований.

Логично, что идею подхода с квадродеревом можно применить и к случаю одномерного массива информации --- аудиоданных. Только вместо деления на 4 сектора можно, например, использовать деление отрезка пополам. То есть суть метода остаётся прежней: имеем разбиение на крупные ранговые блоки, а также пороговое значение $\varepsilon$. В случае, если не удаётся добиться достаточно хорошего сходства для некоторого $R_i$, прибегаем к более мелкому делению.

Интересным способом выделения блоков в изображении также является разбиение на треугольники. Более подробно о нём можно почитать в статье В.\,Г.\,Колебошина и А.\,В.\,Крапивенко, где проводится сравнение этого метода с более традиционными \cite{triangularDiv}.

\subsection{Метрика}

Важным вопросом, на который нужно уметь ответить для реализации алгоритма, является определение степени схожести двух фрагментов изображения (звука). Эта тема является не такой простой, как может показаться на первый взгляд, поскольку восприятие человеком визуальной и слуховой информации --- это комплексный процесс.

Для начала, выдвинем несколько требований к метрике в контексте нашего алгоритма.
\begin{itemize}
  \item Метрика должна в целом отражать сходство данных. То есть близкие к нулю значения расстояния между изображениями (аудио) должны сигнализировать о их сходстве.
  \item Метрика должна быть не слишком ресурсозатратна для вычисления, поскольку она будет посчиттана огромное количество раз в ходе работы программы.
\end{itemize}

Какие меры расстояния можно расмотреть в качестве кандидатов на роль метрики в нашем алгоритме? В основном, выбор был между следующими кандидатами.
\begin{enumerate}
  \item Расстояние Чебышёва: $\rho(x,y) = \max\limits_{i}|x_{i}-y_{i}|$.
  \item Евклидово расстояние: $\rho(x,y) = \sqrt{\sum\limits_{i = 1}^{n}(x_{i}-y{i})^2}$.\\Часто используется его модификация Root Mean Square Error (RMSE): \\$\rho(x,y) = \sqrt{\frac{1}{n}\sum\limits_{i = 1}^{n}(x_{i}-y{i})^2}$.
  \item Расстояние Минковского: $\rho(x,y) = (\sum\limits_{i = 1}^{n}|x_{i}-y{i}|^p)^{1/p}$.
  \item Peak Noise to Signal Ratio (PSNR): $\rho(x,y) = 20\log_{10}(\dfrac{MAX_I}{RMSE})$, где $MAX_I$ --- максимальное значение, принимаемое пикселем изображения или образцом звуковых данных. Для 8 бит это 255.
\end{enumerate}

Если рассматривать метрики с точки зрения схожести изображений, то наиболее близкие к человеческому восприятию результаты показывает PSNR. Хотя, конечно, любые метрики могут давать результаты, сильно расходящиеся с нашим представлением о похожих фрагментах. Так расстояние Чебышёва, как нетрудно догадаться, при изменении всего одного пикселя изображения будет сигнализировать о сильных расхождениях, а Евклидово расстояние отреагирует на смену яркости и контраста, которые не сильно заметны глазу. Для звука справедливы все вышеупомянутые утверждения.

Что же касается вычислительной сложности, очевидно, что проще всего вычислять расстояние Чебышёва, но оно не всегда даёт желаемые нами результаты. Расстояние Минковского со степенями больше двух требует довольно больших вычислений, поэтому его использование мы посчитали нерациональным. Остался выбор между RMSE(так как эта вариация евклидова расстояния более уместна в рамках настоящей работы) и PSNR. В программе реализованы обе, но по умолчанию используется RMSE, так как вычисляется она проще, а различия получившихся в ходе декомпрессии результытов с этими метриками минимальны.

\subsection{Преобразования}

Одно из главных отличий алгоритмов фрактального сжатия для разных предметных областей --- это допустимые преобразования, применяемые к блокам данных. В контексте изображений речь идёт о нижеперечисленных.
\begin{itemize}
  \item Сжатие. Поскольку доменные блоки, как правило, большего размера, чем ранговые, необходимо уменьшить их программно до размера ранговых.
  \item Поворот. Для достижения наилучшего сходства необходимо уметь поворачивать блоки.
  \item Отражения. В комбинации с поворотами позволяют задать нужные для хорошего соответствия преобразования.
  \item Цветокоррекция. Одним из важнейших этапов является корректировка цвета пикселей блоков.
\end{itemize}

Проблема алгоритма цветокоррекции решается благодаря выбранной нами метрике. По сути алгоритм нахождения коэфициентов изменения контрастности и яркости выводится из минимизации метрики RMS.

В алгоритме сжатия изображений реализованы 4 поворота: 0, 90, 180 и 270 градусов; одно отражение по горизонтали или его отсутствие. Также реализован алгоритм сжатия блока в число раз, являющееся степенью 2. Все эти преобразования комбинируются во время перебора для достижения наилучшего сходства рангового и изменённого доменного блока.

Аудиоинформация может быть подвержена следующим преобразованиям.
\begin{itemize}
  \item Сжатие. Оно тут служит той же цели, что и в случае изображений --- изменение доменных блоков до размера ранговых.
  \item Коррекция звука при помощи коэфициентов, вычисленных из метрики. Считается аналогично цветокоррекции изображений.
\end{itemize}

Алгоритм сжатия аудиоданных имеет заметно меньше преобразований, что логично, ибо звук --- одномерный массив данных, а изображение --- двумерный, так что для изображения возможны дополнительные пространственные модификации (повороты и отражения). Интересно, что наиболее важным для сходимости является применение вычисленных из метрики коэфициентов. Это преобразование отражает волновую природу аудиосигнала и цвета, выполняя изменение частоты волны.

\subsection{Метод перебора}

Поскольку алгоритм подразумевает большое число сравнений блоков и перебора параметров преобразований, очень важно правильно организовать сам переборный механизм.

Первоначальной идеей было для каждого рангового блока перебирать все доменные блоки, каждый раз применяя к ним все преобразования. Но этот подход оказался крайне непрактичным и приводил к очень большим затратам времени из-за постоянного пересчитывания одних и тех же величин раз за разом. Ведь даже в случае аудиоданных были лишние изменения масштаба при каждой повторной встрече доменного блока.

После логика перебора была переписана полностью, теперь алгоритм сначала генерировал всевозможные вариации доменных блоков, для которых не нужно знать ничего о ранговых: то есть каждый доменный блок изображения сжимался, всячески отражался и поворачивался, в итоге получалось на каждый исходный домен по 8 вариантов (4 поворота $\times$ 2 отражения) для изображений и по одному варианту для звука (что всё ещё полезно для избежания вышеуказанных проблем с производительностью). Каждый получившийся сжатый блок, а также информация о его положении и применённых преобразованиях хранится внутри структуры, что позволяет избавиться от лишних действий в дальнейшем. Далее, уже при переборе для каждого рангового блока к доменам применялось только изменение цвета, поскольку оно зависело от обоих фрагментов изображения. Подобная модификация в разы увеличила скорость работы алгоритма.

Пример структуры, хранящей информацию о вариации доменного блока изображения:
\begin{verbatim}
type DomainVariant struct {
    X        int
    Y        int
    Rotation int
    Flip     int
    Block    [][]uint8
}
\end{verbatim}

Аналогичная структура, описывающая доменный блок аудиоданных:
\begin{verbatim}
type DomVariant struct {
    Pos   int
    Block []int
}
\end{verbatim}

\subsection{Псевдокод алгоритма}

После того, как мы рассмотрели составные части алгоритма компрессии, мы можем описать более предметно, как происходит процесс сжатия изображений (алгоритм 1).

\begin{algorithm}
  \floatname{algorithm}{Алгоритм}
  \algrenewcommand\algorithmicrequire{\textbf{Вход: }}
  \algrenewcommand\algorithmicensure{\textbf{Выход: }}
  \caption{Псевдокод основных этапов сжатия изображений}\label{alg1}
  \begin{algorithmic}[1]
    \Require Изображение в формате PNG или JPEG
    \Ensure Файл, содержащий коэфициенты преобразований
    \State Считать изображение и конвертировать его в оттенки серого
    \State Добавить рамку вокруг изображения, если его длина или ширина не деляться нацело на размер рангового блока.
    \State Разбить изображение на доменные блоки, каждый из них сжать до размеров ранговых.
    \State Сгенерировать всевозможные вариации доменных блоков, комбинируя повороты и отражения.
    \For{$R_i \in$ мн-во ранговых блоков изображения}
      \For{$Dv_j \in$ мн-во вариаций доменных блоков}
        \State Вычислить коэфициенты контрастности и яркости, дающие минимальную метрику для данных блоков.
        \State Применить вычисленные коэфициенты.
        \State Вычислить метрику RMS.
        \State Если метрика лучше прошлой, запомнить коэфициенты преобразования.
      \EndFor
      \State Запомнить коэфициенты наилучшего преобразования для $R_i$
    \EndFor
    \State Записать коэфициенты всех наилучших преобразований для всех блоков в файл
  \end{algorithmic}
\end{algorithm}

Помимо сжатия изображений мы также можем подробнее рассмотреть процесс сжатия звука (алгоритм 2). Стоит сделать важное уточнение, что алгоритм поддерживает стереозвук, поэтому сжимает оба канала аудио. Делается это просто запуском алгоритма сжатия для каждого канала в отдельности.
\begin{algorithm}
  \floatname{algorithm}{Алгоритм}
  \algrenewcommand\algorithmicrequire{\textbf{Вход: }}
  \algrenewcommand\algorithmicensure{\textbf{Выход: }}
  \caption{Псевдокод основных этапов сжатия звука}\label{alg2}
  \begin{algorithmic}[1]
    \Require Аудиоданные в формате WAV
    \Ensure Файл, содержащий коэфициенты преобразований
    \State Считать звук из файла и разбить на 2 канала: L и R.
    \State Добавить в конец записи несколько значений, чтобы привести длину к кратной размеру рангового блока.
    \State Разбить данные на доменные блоки, каждый из них сжать до размеров ранговых.
    \For{$R_i \in$ мн-во ранговых блоков изображения}
      \For{$Dv_j \in$ мн-во вариаций доменных блоков}
        \State Вычислить коэфициенты, дающие минимальную метрику для данных блоков.
        \State Применить вычисленные коэфициенты.
        \State Вычислить метрику RMS.
        \State Если метрика лучше прошлой, запомнить коэфициенты преобразования.
      \EndFor
      \State Запомнить коэфициенты наилучшего преобразования для $R_i$
    \EndFor
    \State Записать коэфициенты всех наилучших преобразований для всех блоков в файл
  \end{algorithmic}
\end{algorithm}

Как можно заметить, перебор фрагментов для каждого рангового блока происходит независимо от других, поэтому данный алгоритм очень легко параллелизируется. В представленной здесь реализации программа пользуется этим хорошим свойством фрактального сжатия и позволяет пользователю назначить любое количество потоков, значительно увеличивая скорость работы.

\subsection{Декодирование}


Поскольку отображение, построенное из совокупности преобразований доменных блоков в ранговые, является сжимающим с неподвижной точкой, близкой к исходному изображению, процесс восстановления картинки по кодам, считанным из файла, получается довольно простым.
\begin{enumerate}
  \item Считываем из файла размер изображения (аудиофайла), доменных и ранговых блоков.
  \item Считываем данные о преобразованиях.
  \item Генерируем белый шум. В целом можно взять любой фрагмент данных нужных размеров.
  \item Итеративно применяем к белому шуму преобразования доменных блоков в ранговые.
  \item Спустя небольшое число итераций получаем изображение (звук).
\end{enumerate}

Поскольку отображение сжимающее и его неподвижная точка является исходным изображением (звуком), не важно, с каких входных данный мы начнём, в результате мы придём к тому, что было закодировано. Пример различных стартовых изображений можно видеть на рисунках \ref{cat_stages} и \ref{cat_stages_dog}.

\begin{figure}
  \begin{center}
    \includegraphics[width=0.95\textwidth]{./images/cat_stages.jpg}
  \end{center}
  \caption{Пример декодирования. Начальное значение --- белый шум}\label{cat_stages}
\end{figure}

\begin{figure}
  \begin{center}
    \includegraphics[width=0.95\textwidth]{./images/cat_stages_dog.jpg}
  \end{center}
  \caption{Пример декодирования. Начальное значение --- другое изображение}\label{cat_stages_dog}
\end{figure}

\subsection{Хранение коэффициентов в файле}

Одной из самых трудных и творческих задач было разработать формат файлов, позволяющий эффективно хранить данные о преобразованиях. В ходе работы над алгоритмом и его отладки в качестве временного решения использовался JSON. Этот формат файлов очень хорошо помогает визуально оценить правильность нахождения соответствий и в целом визуализировать работу программы, однако не годится в качестве метода хранения сжатого файла, поскольку является текстовым файлом, ещё и с дополнительными символами, такими как ковычки или скобки.

По причине необходимости компактного хранения данных было решено разработать свой формат файлов. Требования, которые были к нему предъявлены.
\begin{itemize}
  \item Использование бинарного файла, а не текстового, так как текстовый занимал бы слишком много места.
  \item Отсутствие лишнего места между сериализованными (записанными в файл) данными. То есть если поле структуры имеет тип 8-битное целое число, но принимает только значения от 0 до 3, то в файле оно должно занимать не более двух бит.
  \item Вариативная длина полей, отвечающих за координаты. Поскольку координаты доменных блоков всё-таки нужно хранить, под их хранение выделяется минимально возможное количество бит.
\end{itemize}

Говоря о сжатом представлении изображения, речь идёт о двух следующих структурах, где Point --- просто структура с координатами x, y:
\begin{verbatim}
type CompressedData struct {
    Width   int
    Height  int
    RngSize int
    DomSize int
    Matches []Match
}

type Match struct {
    Rng         Point
    Dom         Point
    Rotation    int
    Flip        int
    ColorScale  float64
    ColorOffset int
}
\end{verbatim}

Информация о сжатом звуке ещё более краткая:
\begin{verbatim}
type Compressed struct {
    LeftMatches  []Match
    RightMatches []Match
    Length       int
    RngSize      int
    DomSize      int
}

type Match struct {
    DPos    int
    RPos    int
    Scale   float64
    Offset  int
}
\end{verbatim}

Можно сделать несколько наблюдений, анализируя работу программы.
\begin{itemize}
  \item Мы точно знаем разбиение на доменные и ранговые блоки, если нам известны длина, ширина изображения (или длина звука) и размеры этих блоков.
  \item Поворотов может быть только 4, отражений только 2 (никакое или по горизонтали).
  \item Разбиение на доменные блоки происходит в позициях, кратных величине рангового блока. Это сделано для того, чтобы перебор был не слишком долгим, но при этом было достаточное количество доменных блоков для построения хорошего соотношения. Это можно использовать для более компактного хранения позиции блоков.
\end{itemize}

Из первого наблюдения следует, что мы можем, например, не хранить вообще позицию рангового блока, вместо этого упорядочив все эти блоки, и в таком порядке написать их коэфициенты в файл. Обратное восстановление также возможно, благодаря известному порядку.

Второе наблюдение указывает на то, что все данные об отражениях и поворотах можно записать в 3 байта.

Наконец, третье позволяет сэкономить на кодировании доменных блоков. Вместо того, чтобы запоминать позицию по x и y в изображении, можно вместо этого запомнить номер блока по ширине и номер блока по высоте. И поскольку эти величины напрямую связаны с размерами изображения и ранговых блоков, получается, что они не тратят лишние байты памяти, когда это не нужно.

Помимо этого было решено ограничиться точностью коэфициента контрастности до трёх знаков после запятой. А так как коэфициент принимает значения в диапазоне [0, 1], его можно хранить в виде десятибитного целого числа (достаточно умножить его на 1000: если $x\in[0,1]$, то $x\times 1000\in[0,1000] < 1024$). Аналогично хранится и дробный коэфициент при сжатии аудио.

Размеры изображения были ограничены 16-битными положительными числами (до 65536 по ширине или высоте). Ограничения вполне оправданы, так как маловероятно, что при такой временной сложности на вход будут переданы настолько большие изображения. Для аудио размер был ограничен 32 битами.

Для записи файла был также написан класс, способный читать и писать биты в файл без пустых мест между ними (разве что последний байт файла будет дополнен нулями, потому что не может быть файлов с нецелыми байтами в памяти).

В итоге получился следующий алгоритм записи файла.
\begin{enumerate}
  \item Кодируем размеры изображения (аудио) и блоков.
  \item Подсчитываем, какой размер будет занимать Match.
  \item Итерируемся по структурам Match и конвертируем каждую в 64-битное целое число.
  \item Зная настоящий размер численного представления структуры, пишем в бинарный файл ровно столько бит, сколько она занимает.
\end{enumerate}

\section{Программная реализация}

Этот раздел соосредоточен на рассмотрении того, как описанный выше алгоритм реализован в виде кода на языке Go. Здесь будет рассмотрено, как соотносятся этапы работы алгоритма с функциями и структурами программы. Полный код разработанных в рамках этой работы программ представлен в приложении А.

\subsection{Сжатие}

Программа для сжатия изображений работает с изображениями в форматах JPEG и PNG, которые конвертируются в оттенки серого перед стартом алгоритма, поскольку тот работает только с чёрно-белыми изображениями. Аудио принимается в формате WAV, причём поддерживается стерео. Для конвертации других аудиофайлов в WAV можно воспользоваться бесплатной программой FFmpeg.

Работа программы начинается с создания структуры \texttt{Compressor} посредством вызова функции \texttt{NewCompressor}. Эта функция принимает на вход названия исходного файла, а также того, в котором стоит сохранить результат, размеры ранговых и доменных блоков, количество потоков, а также пороговое значение для перебора. Внутри этой функции считываются данные из файла, который необходимо сжать (в случае изображений чтение выполнено с помощью пакета image стандартной библиотеки Go, в то время как звук считывается при помощи пакета go-wav), после этого данные для кодирования представляются в виде массива целых чисел в структуре \texttt{Compressor}.

У структуры \texttt{Compressor} определён метод \texttt{Compress}, в котором происходит основная работа алгоритма. Сначала создаётся объект типа \texttt{Comp\-res\-sed\-Da\-ta}, заполняются размеры изображения (аудиофайла), ранговых и доменных блоков. После этого при помощи функции \texttt{generateDomVariants} происходит заполнение поля \texttt{DomVariants} структуры \texttt{Compressor} всеми возможными вариациями доменных блоков (после применения преобразований).

Для заполнения информации об отображениях доменных блоков в ранговые, вызывается функция \texttt{runParallel}. Эта функция запускает перебор в указанном при создании компрессора числе потоков для эффективного использования вычислительных мощностей компьютера, а также в дополнительном потоке работает функция, отвечающая за печать прогресса в консоль.

Сам перебор происходит при помощи двух функций: \texttt{findBestMatches} и \texttt{findBestMatch}. Вторая отвечает за просмотр всех доменных блоков и поиск наилучшего совпадения для заданного рангового, в то время как первая вызывает вторую функцию для всех ранговых блоков в определённом диапазоне (нужно это для того, чтобы можно было в каждом потоке вызвать \texttt{findBestMatches} для своего кусочка исходных данных).

После того, как \texttt{runParallel} завершит свою работу, происходит запись результатов работы программы в сжатом представлении в файл.

\subsection{Декодирование}

Процесс декодирования заметно короче. Начинается всё также с создания структуры \texttt{Decompressor} при помощи функции \texttt{NewDecompressor}. В этой функции происходит чтение сжатого файла, из которого восстанавливается структура \texttt{CompressedData}. Информация из \texttt{CompressedData} копируется в поля \texttt{Decompressor}, помимо этого там также хранится информация о степени масштабирования итогового результата и фрагмент данных, к которому будут итеративно применяться преобразования, считанные из файла (функция \texttt{NewDecompressor} инициализирует значения этого фрагмента случайными значениями).

У структуры \texttt{Decompressor} определён метод \texttt{Decompress}, который реализует сам процесс декодирования. В нём итеративно применяются преобразования из файла. После завершения декодирования, итоговый результат сохраняется опять же при помощи функций из библиотек image и go-wav.


\subsection{Консольный интерфейс}

Помимо всего прочего для программы был разработан консольный интерфейс. Он представляет собой набор команд и флагов, позволяющий удобно управлять параметрами алгоритма для достижения желаемых результатов. Написание CLI-интерфейса выполнено при помощи библиотеки go-flags \cite{goflags}, позволяющей автоматизировать парсинг аргументов командной строки и гибко настроить поддерживаемые флаги и команды при помощи тегов, приписываемых полям структур. Также приятным бонусом является автоматическая генерация сообщения с инструкцией по использованию при введении ключа \texttt{help}.

Консольный интерфейс позволяет настраивать в сжатии параметры, указанные ниже.
\begin{itemize}
  \item Название входного и выходного файла.
  \item Размеры ранговых и доменных блоков.
  \item Количество потоков для перебора.
  \item Пороговое значение для оптимизации перебора и преждевременой остановки в случае нахождения достаточно хорошего совпадения.
\end{itemize}
В декодировании, в свою очередь, можно указать следующие параметры.
\begin{itemize}
  \item Название входного и выходного файла.
  \item Масштабирование изображения при декомпрессии.
\end{itemize}

Консольный интерфейс одинаков для обеих программ, поскольку параметры, настраиваемые в компрессии аудио и изображений, идентичны. На рисунке \ref{interface} представлен пример того, как этот интерфейс выглядит, полностью его можно увидеть в приложении B.

\begin{figure}
  \begin{center}
    \includegraphics[width=0.95\textwidth]{./images/cli-example.png}
  \end{center}
  \caption{Пример справки консольного интерфейса}\label{interface}
\end{figure}

Помимо этого, реализован счётчик прогресса многопоточного перебора, позволяющий видеть, какой процент соответствий найден суммарно всеми потоками. Эта проблемма, которая бы потребовала в C/C++ использования, например сокетов, изящно решается при помощи особого примитива синхронизации в Go --- каналов, которые являются удобным средством коммуникации между разными потоками одного приложения.

\section{Фрактальное сжатие  других типов информации}

Как было видно из сходства двух написанных программ, а также из описания общей идеи алгоритма, фрактальное сжатие концептуально может быть применено к любому типу информации. Для того, чтобы создать алгоритм сжатия для какой-то конкретной предметной области, нужно ответить на следующие вопросы.
\begin{enumerate}
  \item Как провести разбиение сжимаемой информации на блоки?
  \item Какие преобразования доменных блоков возможны?
  \item Какая метрика используется для определения степени сходства?
  \item Как хранятся сжатые данные?
  \item Как декодировать файл после сжатия?
  \item Каким образом применять коэфициенты преобразований к блокам при декодировании?
\end{enumerate}

Поскольку все эти вопросы довольно специфичны для сферы применения алгоритма, написать какой-то универсальный алгоритм не получится. Однако, придерживаясь общей схемы и внося изменения, диктуемые предметной областью, можно представить алгоритм для этой области.

Также стоит помнить о серьёзном недостатке фрактального способа компрессии --- долгом времени сжатия. Этот недостаток делает алгоритм неприменимым во многих предметных областях. Поскольку изначальной сферой применения фрактального сжатия были изображения, после чего были исследования, связанные с сжатием звука, логичным развитием является фрактальное сжатие видеоинформации. Исследования по этой теме ведутся и по сей день \cite{fractVideo1, fractVideo2}. 

Сжатие видео должно при создании соотношений работать не в рамках одного кадра, сжимая каждое изображение по отдельности, а в рамках отрезка видео, иначе не добиться хороших показателей компрессии, однако скорость работы алгоритма в таком случае становится главной проблемой, не позволяющей применять его в практических задачах.

Тем не менее, алгоритм действительно применим к различным типам информации, если ответить на вышеописанные вопросы в случае конкретной предметной области.

\section{Результаты работы программы}

Рассмотрим работу программ, реализованных в рамках данной работы на конкретных примерах. В этом разделе будут представлены результаты алгоритмов для изображений и звука, сравнения размера, занимаемого сжатым представлением и оригиналом, а также качества восстановленных данных в сравнении с исходными.

\subsection{Сжатие изображений}

В качестве тестовых были выбраны изображения, представленные на рисунке \ref{testimg}. В ходе эксперимента мы сожмём их, оценим отличия в качестве и размере файлов.
\begin{figure}
    \includegraphics[width=0.5\textwidth]{./images/lena.png}
    \includegraphics[width=0.5\textwidth]{./images/cat.png}
    \caption{Тестовые изображения}\label{testimg}
\end{figure}

Начнём с того, что приведём результаты промежуточных шагов декодирования для иллюстрации того, как при помощи сжимающего отображения мы итеративно приближаемся к изображению, что было закодировано. На рисунке \ref{manystages} показан процесс восстановления тестовых изображений.


\begin{figure}
  \begin{center}
    \includegraphics[width=0.95\textwidth]{./images/lena_stages.jpg}
  \end{center}
  \begin{center}
    \includegraphics[width=0.95\textwidth]{./images/cat_stages.jpg}
  \end{center}
  \caption{Шаги декодирования}\label{manystages}
\end{figure}


Рассмотрим, как именно можно проверить качество восстановленного изображения. Для этого можно использовать уже знакомую нам метрику PSNR: чем она ниже, тем лучше сходство. Далее представлена таблица, в которой для каждого из изображений при разных размерах доменных и ранговых областей подсчитана метрика PSNR, а также указан коэфициент сжатия (отношение размера исходного несжатого файла к размеру сжатого представления). В качестве несжатоно файла используется изображение в формате BMP.

\begin{center} % cat.bmp = 426KB  lena.bmp = 362KB
\begin{tabular}{|c|c|c|c|}
  \hline
  Изображение, параметры & PSNR & Коэфициент сжатия & Время работы\\\hline
  cat 736$\times$736, r2d4 & 38dB & 0.75 & 35.4 с\\
  cat 736$\times$736, r4d8 & 42dB& 2.7 & 31.1 с \\
  cat 736$\times$736, r4d16 & 41dB & 2.7 & 18.7 с \\
  cat 736$\times$736, r8d16 & 43 dB & 11.5 & 11.6 с \\
  lena 512$\times$512, r2d4 & 30dB & 1.13 & 11.5 с \\
  lena 512$\times$512, r4d8 & 38dB & 4.76 & 9.8 с \\
  lena 512$\times$512, r4d16 & 35dB & 4.76 & 5.1 с \\
  lena 512$\times$512, r8d16 & 42dB & 20 & 3.2 с \\\hline
\end{tabular}
\end{center}

Как видно из экспериментальных данных, при размере рангового блока 2$\times$2 пикселя уменьшение объёма происходит не всегда, но вот при более крупных ранговых блоках коэфициент сжатия стремительно увеличивается, однако при этом также сильно страдает и качество картинки. В целом показатели как качества, так и степени сжатия далеки от максимально хороших и могут быть заметно улучшены, если применить другой подход к разбиению (квадродерево), а также доработать представление в виде файла.

Примеры изображений при разных параметрах алгоритма изображены ниже. Можно заметить, как при увеличении размеров рангового блока качество восстановленного изображения заметно падает, но при этом вырастает степень сжатия (рисунки \ref{imgr2d4}--\ref{imgr5d10}).

\begin{figure}[H]
  \includegraphics[width=0.5\textwidth]{./images/catr2.jpg}
  \includegraphics[width=0.5\textwidth]{./images/lenar2.jpg}
  \caption{Ранговый блок $2\times 2$ пикселя, доменный $4\times 4$}\label{imgr2d4}
\end{figure}

\begin{figure}[H]
  \includegraphics[width=0.5\textwidth]{./images/catr4.jpg}
  \includegraphics[width=0.5\textwidth]{./images/lenar4.jpg}
  \caption{Ранговый блок $4\times 4$ пикселя, доменный $8\times 8$}
\end{figure}

\begin{figure}[H]
  \includegraphics[width=0.5\textwidth]{./images/catr5.jpg}
  \includegraphics[width=0.5\textwidth]{./images/lenar5.jpg}
  \caption{Ранговый блок $5\times 5$ пикселей, доменный $10\times 10$}\label{imgr5d10}
\end{figure}

\subsection{Сжатие аудио}

Экспериментальные исследования производились с двумя аудиофайлами: отрывком из <<Лунной Сонаты>> Людвига Ван Бетховена (MS.wav), а также фрагментом трека <<Pale Blue Eyes>> группы The Velvet Underground (PBE.wav). Оба аудиофайла обрезаны до 1 минуты, а также закодированы в формат WAV (2 канала, 8000 Гц, 128 кбит/с). Проводились такие же исследования, как и в случае с изображениями.

\begin{figure}[H]
  \begin{center}
  \includegraphics[width=0.95\textwidth]{./images/MS-init.jpg}
  \end{center}
  \caption{<<Лунная Соната>>, левый аудиоканал}
\end{figure}

\begin{figure}[H]
  \begin{center}
  \includegraphics[width=0.95\textwidth]{./images/PBE-init.jpg}
  \end{center}
  \caption{<<Pale Blue Eyes>>, левый аудиоканал}
\end{figure}

На рисунке ниже можно увидеть, как изменяются файлы в ходе работы алгоритма декодирования, постепенно приближаясь к исходному сжатому звуку. В целях экономии места для иллюстраций на рисунках изображён только левый канал.

\begin{figure}[H]
  \begin{center}
  \includegraphics[width=0.95\textwidth]{./images/ms_stages.jpg}
  \end{center}
  \caption{Декодирование <<Лунной Сонаты>>}
\end{figure}

Для аудиоданных также была использована метрика PSNR для нахождения степени расхождений между начальным и итоговым аудиофрагментом. Вес кодируемых айдиофайлов составляет 938 килобайт.

\begin{center}
\begin{tabular}{|c|c|c|c|}
  \hline
  Аудио, параметры & PSNR & Коэфициент сжатия & Время работы\\\hline
  MS.wav, r4d8   & 32.4dB & 0.98 & 14 с\\
  MS.wav, r6d18  & 33.8dB & 1.47 & 45 с \\
  MS.wav, r10d20 & 35.6dB & 2.36 & 34 с \\
  MS.wav, r10d30 & 34dB   & 2.36 & 33.6 с \\
  PBE.wav, r2d4  & 31.3dB & 0.98 & 15 с \\
  PBE.wav, r4d8  & 34dB   & 1.47 & 47 с \\
  PBE.wav, r4d16 & 34.8dB & 2.36 & 32 с \\
  PBE.wav, r8d16 & 34.3dB & 2.36 & 33.5 с \\\hline
\end{tabular}
\end{center}

Эксперимент показал, что алгоритм действительно приводит к уменьшению занимаемого файлами места на диске при размере рангового блока больше четырёх. Показатель PSNR отражает утверждение о том, что чем больше ранговые блоки мы берём, тем более зашумлённым получается результат. Однако, в случае с аудиоданными падение качества не такое заметное, как в случае изображений.

В целом, рассмотрев примеры работы программы по сжатию изображений и аудио, можно сделать вывод, что в текущей реализации программы хоть и неконкурентноспособны с популярными алгоритмами сжатия для своих предметных областей (для конкурентных результатов требуется заметно усложнить структуру программ и добавить различные эвристики), однако они хорошо служат в качестве демонстрации возможностей алгоритма и показывают, как сжимающие отображения и фракталы могут быть использованы для компактного хранения данных.

\phantomsection
\addcontentsline{toc}{chapter}{Заключение}
\chapter*{Заключение}

В ходе работы были изучены основные алгоритмы фрактального сжатия изображений и звука, их модификации, а также рассмотрена теория, стоящая за этими алгоритмами.

На основе изученных алгоритмов реализованы программы сжатия данных. Программы реализуют классический метод фрактального сжатия при помощи PIFS с разбиением на блоки фиксированного размера. Это довольно прямолинейный подход, плюсы и минусы которого коррелируются с полученными на практике результатами. В работе продробно расписаны составные части программ и потенциальные пути для модиикации, например более эффективное деление на блоки различных размеров для уменьшения итогового объёма файлов при сохранении приемлемого качества.

Для программ разработан удобный интерфейс и формат сжатого представления данных в файле. Файл является ещё одним потенциальным местом для модификаций для улучшения показателей сжатия.

Кроме того, в работе освещён вопрос применимости фрактального кодирования к другим типем информации. Поскольку алгоритм в своём абстрактном виде применим к практически любым предметным областям, его использование признано возможным --- нужно только определить преобразования, метрики, разбиения и т.\,п. для конкретных данных. Однако стоит помнить о больших вычислительных затратах, способных стать критическими во многих задачах.

\printthebibliography

\appendix
\chapter{Код программ фрактального сжатия}
\section{Программа сжатия изображений}
\chapter{Код программ фрактального сжатия}
\section{Программа сжатия изображений}
\chapter{Код программ фрактального сжатия}
\section{Программа сжатия изображений}

\end{document}
